\section{Design}
\label{sec:design}

% Decision pipeline for one proposed tool call.
\begin{figure*}[t]
\centering
\begin{tikzpicture}[
  node distance=3.5mm and 2.2mm, font=\scriptsize,
  box/.style={draw, rounded corners=2pt, align=center, minimum height=9mm, inner sep=2pt, fill=white, text width=24mm},
  layer/.style={box, fill=cA!7, draw=cA!60},
  state/.style={box, fill=cC!8, draw=cC!60, dashed, text width=40mm},
  outbox/.style={box, fill=cB!8, draw=cB!60, text width=28mm},
  arr/.style={-{Stealth[length=1.6mm]}, thick}]
\node[box] (host) {Host hook\\(Claude Code, OpenHands,\\SDK, MCP proxy)};
\node[layer, right=of host] (l0) {L0\\canonicalise\\(NFKC, hidden chars)};
\node[layer, right=of l0] (l1) {L1 policy\\tiers (+\,\texttt{when}), scope,\\argument rules, clamps,\\redaction, budgets};
\node[layer, right=of l1] (ty) {Typed actions\\targets, ports,\\rates as values};
\node[layer, right=of ty] (l2) {L2 predictive\\effect on the\\asset inventory};
\node[layer, right=of l2] (fl) {Session flow\\untrusted / sensitive\\/ egress};
\node[outbox, below=6mm of fl] (v) {worst finding wins;\\allow + repair $\to$ \verdict{modify};\\escalate after $N$ denials};
\node[outbox, left=of v] (ap) {\verdict{ask}: approval queue\\(reuse, single use, TTL);\\calibrated deferral\\(tier-only asks)};
\node[outbox, left=of ap] (led) {hash-chained ledger;\\runtime monitor};
\node[state, below=7mm of host.south west, anchor=north west] (st) {session state: denials, budgets used,\\assets taken down, data labels\\(read by L1, L2 and flow)};
\draw[arr] (host) -- (l0);
\draw[arr] (l0) -- (l1);
\draw[arr] (l1) -- (ty);
\draw[arr] (ty) -- (l2);
\draw[arr] (l2) -- (fl);
\draw[arr] (fl) -- (v);
\draw[arr] (v) -- (ap);
\draw[arr] (ap) -- (led);
\draw[arr, dashed] (st.east) -| (l1.south);
\draw[arr, cB] (led.south) |- node[pos=0.75, below] {verdict to the host; any error $\to$ \verdict{deny}} ([xshift=-3mm,yshift=-3mm]st.south west) |- (host.west);
\end{tikzpicture}
\caption{\SystemNameCap{} judges one proposed call. Every layer adds findings; the strictest finding decides. Calls that are let through (including approved and auto-approved ones) update the session state that later calls are judged against.}
\label{fig:architecture}
\end{figure*}


\subsection{Decision contract}
A host sends a proposed call: a session identifier, a tool name and JSON arguments, optionally an agent and framework name. The answer is one of four verdicts with a risk level, the findings (rule, verdict, reason) that produced it, the corrected arguments for \verdict{modify}, an approval identifier for \verdict{ask}, and the reports of the predictive, flow and typed layers. Each layer adds findings; the verdict is the most severe ($\verdict{allow} < \verdict{modify} < \verdict{ask} < \verdict{deny}$), and an \verdict{allow} whose arguments were repaired becomes \verdict{modify} (Fig.~\ref{fig:architecture}).

\begin{property}[Fail closed]
An exception anywhere in the decision, a timeout (2\,s by default), arguments above a size limit, a second concurrent check for the same session, a malformed hook payload, and a missing, pending, expired, used or mismatched approval each produce \verdict{deny}. A call identical to one already denied in the session stays denied until a person resets the session, and repeats count towards escalation: after $N$ denials (three by default) every call that would be allowed needs approval.
\end{property}

\subsection{Canonicalisation and policy tiers}
Before any rule runs, tool names and argument strings are normalised (Unicode NFKC) and stripped of invisible, bidirectional and tag characters; a hidden character in a tool name is a \verdict{deny}, in an argument an \verdict{ask}, because what a reviewer would see is not what runs. Fields that carry the agent's stated rationale (\texttt{reasoning}, \texttt{justification} and similar) are ignored by every rule (G3).

A policy pack maps tool names (globs) to \emph{tiers}, each with a decision and a reversibility flag, for example \emph{read} $\to$ \verdict{allow}, \emph{reversible containment} $\to$ \verdict{ask}, \emph{irreversible} $\to$ \verdict{deny}; unknown tools get the pack's default tier. A tool rule can be limited to calls whose arguments fully match regular expressions (\texttt{when}), which expresses ``allow \texttt{git status} but ask for any other shell command''; conditions may not test rationale fields. Further rules add findings: \emph{scope} (targets must lie inside agreed networks and host globs; with \emph{projection}, an out-of-scope list entry is dropped or a broad range narrowed to the single allowed network inside it, and the call is denied only when nothing in scope remains), \emph{argument rules} (regular expressions on argument paths, e.g.\ shell chaining, credential paths, agent-configuration files), \emph{clamps} (numeric caps become \verdict{modify}), \emph{redaction} (secret-bearing arguments are removed) and \emph{budgets} (per-session totals of calls or of a numeric argument, with a reserve; amounts over what is left are clipped, exhausted budgets stop the call). A call that needs approval carries its budget charge; the charge is applied when the approved call is used, after checking the budget again.

\subsection{Typed actions}
String rules cannot count the addresses in \texttt{10.0.0.0/1}, the ports in \texttt{1-65535}, or recognise \texttt{0xa9.0xfe.0xa9.0xfe} and \texttt{[::ffff:169.254.169.254]} as the instance-metadata address. For tools a pack declares as network actions, \SystemName{} parses the named target, port and rate arguments into values first: addresses and networks, ranges (\texttt{10.0.0.1-50}), trailing wildcard octets, hostnames (IDNA, trailing dot), \texttt{host:port}, URLs (scheme, the real host after any user information, port) and lists. It then denies forbidden networks and hosts in any spelling, denies non-canonical numeric IP literals (hex, octal, integer or shortened forms that some resolvers accept and filters miss), unsupported URL schemes and backslashes or percent signs in URL hosts (parser differentials~\cite{orangetsai}), checks the scope on networks and cuts lists and ranges to their exact in-scope part, bounds addresses, ports and probes per call, and lowers rates and parallelism to their limits. A value that cannot be parsed is denied, or, where existing integrations name targets inconsistently (SOC containment), sent to a person.

\subsection{Predictive layer}
For containment actions the predictive layer asks what the call would do to the defender's environment. The tenant maintains an inventory of assets (hosts, accounts, IPs, domains, services) with roles, dependents and redundancy groups, each with a minimum number of members that must stay available. An \emph{effect} maps tools to what they make unavailable. A call is unsafe if, counting what the session has already taken down (approved calls included), a redundancy group would fall below its minimum, or the dependents cut off by the call or the session would exceed limits; unsafe predictions are denied, and targets missing from the inventory go to a person. Effects matched by process name or path go to a person rather than being denied, because an attacker could name a process after a critical service to make the guard refuse to stop it. The session is the horizon, in the spirit of a predictive safety filter~\cite{wabersich}: a call is judged on the state it leaves.

\subsection{Session information flow}
Packs label tools \emph{untrusted} (reads content an outsider can influence), \emph{sensitive} (reads private data) or \emph{egress} (sends data out of the organisation), by tool glob and optionally by an argument pattern. A session collects the labels of calls that were let through. A flow rule fires on a call carrying its sink label when the session, counting this call, holds every required label; the default is the Rule of Two~\cite{ruleoftwo}: an egress call after the session read untrusted input and sensitive data goes to a person, with the call that introduced each label named in the reason. Egress is exempt for destinations on an internal list; every destination must be internal, lists are split, nested keys are searched, and a call whose destination cannot be found counts as external. A tool nothing classifies is assumed to carry every label. Labels are coarse (per session, from which tools ran, never from their results or the agent's text), which over-approximates in the safe direction.

\subsection{Approvals and calibrated deferral}
An \verdict{ask} creates an approval request bound to the session and the canonical call; the reviewer sees the exact (possibly repaired) arguments, targets, tier, predicted effect, data labels and budget charge, and the approved call runs with exactly those arguments, once. Because hook processes are stateless, a repeat of the same call in the same session reuses its pending request, and once approved, the next identical call uses the approval (within the approval time-to-live). A denial makes the call sticky-denied.

Calibrated deferral learns which requests a person almost always approves. Each request gets a transparent risk score (logarithm of the predicted blast radius, unknown targets, an unclassified tool, earlier denials in the session). From the tenant's decided requests (approved $=0$, denied $=1$), LTT~\cite{ltt} tests thresholds on a grid of score deciles from the lowest upward with exact binomial tests of $H_0$: ``the denial share among requests scoring at most $t$ is at least $\alpha$'', accepting each threshold whose test rejects at level $\delta$ and stopping at the first that does not. With probability at least $1-\delta$, at most $\alpha$ of future requests under the accepted threshold are ones a person would have denied, provided future requests resemble past ones. This is a statement about reviewers, not about attackers. We therefore make only requests raised by the action tier alone eligible: never refusals, hidden characters, irreversible or unclassified tiers, unknown or unsafe predicted effects, information-flow rules, typed-action limits, out-of-scope targets, exhausted budgets, or escalated sessions. Deferral runs in \emph{shadow} mode (record what would have been auto-approved) before \emph{auto} mode, and only a signed-in person can switch it.

\subsection{Evidence and verifiable evaluation}
Non-allow decisions, approvals, registry pins, session resets, evaluations, policy proposals and deferral settings are appended to a per-tenant hash chain whose integrity anyone can recompute. Evaluations run a labelled suite (cases with an expected verdict or an accepted set, optional session order, invariance groups and expected rule) and report a digest over suite, policy, filter version and verdicts, so a result can be reproduced exactly; comparisons add single-layer ablations, trivial baselines and bootstrap intervals~\cite{efron} over cases or whole sessions.
